\section{Branch Codes and Pruning Constraints}
\label{sec:subproblems}

We now define the branch-code space for the whole ordered instance and isolate the branch decisions that can affect pruning constraints.

Let the vertices be ordered as $v_1,\ldots,v_n$, and write $U_i$ for the predecessor set of $v_i$. When a vertex $q=v_i$ appears as a subscript, $U_q$ means $U_i$; this preserves the paper's vertex notation while making the order explicit.

Let $V_0=\{v_1,\ldots,v_K\}$ be the fixed initial clique in $\mathbb{R}^K$. For $i\le K$, set $U_i=\varnothing$. For every non-seed vertex $i > K$, the predecessor set $U_i \subseteq \{v_1,\ldots,v_{i-1}\}$ has size $K$. The predecessor relation defines a directed acyclic graph with an arc $v_j \to v_i$ whenever $v_j \in U_i$.

Related order-dependent reference structures have been studied through $K$-discretization and $K$-incident graphs~\cite{abud2018k}; here we fix one predecessor system and derive branch shifts from its dependency cones.

\begin{definition}[Fixing set]
\label{def:fixing}
The fixing set of a vertex $v_i$ is
\[
\operatorname{Fix}_U(v_i)
=
\{v_i\}\cup \bigcup_{u\in U_i}\operatorname{Fix}_U(u).
\]
For seed vertices this gives the base case $\operatorname{Fix}_U(v_i)=\{v_i\}$ because $U_i=\varnothing$.
\end{definition}

\begin{definition}[Dependency cone]
\label{def:cone}
The dependency cone of a vertex $v_i$ is
\[
\operatorname{Cone}_U(v_i)
=
\{v_i\}\cup \bigcup_{w:\,v_i\in U_w}\operatorname{Cone}_U(w),
\]
where $w$ ranges over vertices. If no later vertex has $v_i$ as a predecessor, this gives the terminal base case $\operatorname{Cone}_U(v_i)=\{v_i\}$.
\end{definition}

Let $E_D$ denote the discretization edges induced by the predecessor relation, and let $E_P = E\setminus E_D$ denote the pruning edges. The full branch-decision set is
\[
B = V\setminus V_0.
\]
Define the pruning-relevant vertex set
\[
L
=
\bigcup_{\{u,v\}\in E_P}
\left(\operatorname{Fix}_U(u)\cup \operatorname{Fix}_U(v)\right),
\]
and the pruning-constrained branch set
\[
B_{\mathrm{c}}=L\setminus V_0.
\]
The remaining branch decisions
\[
B_{\mathrm{f}}=B\setminus B_{\mathrm{c}},
\qquad
f=\lvert B_{\mathrm{f}}\rvert,
\]
are free with respect to pruning constraints: changing them preserves all discretization edges by construction and cannot change the endpoints of any pruning edge.

The constrained edge set is
\[
F = E_D[L] \cup E_P,
\qquad
E_D[L]=\{\{u,v\}\in E_D \mid u,v\in L\}.
\]
When $E_P=\varnothing$, we have $L=\varnothing$, $B_{\mathrm{c}}=\varnothing$, and all branch decisions are free.

Let $x(s)=(x_a(s))_{a\in V}$ denote the embedding determined by the full branch code $s\in\mathbb{F}_2^B$. The full solution code is
\[
\Xi
=
\{s \in \mathbb{F}_2^B \mid
\|x_a(s)-x_b(s)\|_2=d_{ab}
\text{ for every } \{a,b\}\in E\}.
\]
The constrained solution code is
\[
\Xi_F
=
\{s_{\mathrm{c}} \in \mathbb{F}_2^{B_{\mathrm{c}}} \mid
\|x_a(s)-x_b(s)\|_2=d_{ab}
\text{ for every } \{a,b\}\in F\},
\]
where $s$ is any full branch code whose restriction to $B_{\mathrm{c}}$ is $s_{\mathrm{c}}$. This is well-defined because branch decisions in $B_{\mathrm{f}}$ do not affect edges in $F$. Hence
\[
\lvert \Xi\rvert = 2^f \lvert \Xi_F\rvert.
\]
Our goal is to compute $\lvert \Xi\rvert$ under generic geometric assumptions without running tree-based branch search.
